% @see : https://coopmaths.fr/alea/?uuid=f1208&id=2N30-flash3&n=2&d=10&qcm=0&cd=1&tip=0&alea=ka4d&uuid=97008&id=2N30-flash2&n=1&d=10&qcm=0&cd=1&tip=0&alea=dpSN&uuid=16c8e&id=2N30-flash1&n=2&d=10&cd=1&tip=0&alea=nY7r&uuid=d309b&id=2N30-4&n=1&d=10&s=false&s2=10&s3=true&s4=true&cd=0&tip=0&alea=q00i&uuid=d0fdc&id=2N33-1&n=1&d=10&s=6&s2=false&cd=1&tip=0&alea=Mdan&uuid=08a2a&id=2N31-flash2&n=2&d=10&qcm=0&cd=1&tip=0&alea=w7p1&uuid=dd7d1&id=2N31-flash1&alea=PvQN&uuid=1894f&id=2N31-flash4&alea=C68e&uuid=07df0&id=2N31-flash3&alea=Ship&uuid=6575c&id=2N32-1&n=1&d=10&s=4&s2=true&cd=0&tip=0&alea=xvAD&uuid=4c878&id=2N33-flash1&alea=LIo0&v=eleve
\documentclass[a4paper,12pt,fleqn]{article}
\usepackage[margin=1.7cm]{geometry}

%%% COULEURS %%%
\usepackage[table,svgnames]{xcolor}
\definecolor{nombres}{cmyk}{0,.8,.95,0}
\definecolor{gestion}{cmyk}{.75,1,.11,.12}
\definecolor{gestionbis}{cmyk}{.75,1,.11,.12}
\definecolor{grandeurs}{cmyk}{.02,.44,1,0}
\definecolor{geo}{cmyk}{.62,.1,0,0}
\definecolor{algo}{cmyk}{.69,.02,.36,0}
\definecolor{correction}{cmyk}{.63,.23,.93,.06}
\arrayrulecolor{couleur_theme} % Couleur des filets des tableaux
\usepackage{tasks}


\usepackage[most]{tcolorbox} % Supprimé à cause de conflits avec ProfCollege
% Découpe des cadres d'exercice. Le style est redéfini par la sortie LaTeX
% juste avant un exercice à garder d'un seul tenant, puis rétabli :
%   \tcbset{mathaleaexo/.style={unbreakable}}
% Passer par un style plutôt que par la clé « breakable » directement évite
% de changer la découpe des autres cadres du document (ProfCollege...).
\tcbuselibrary{documentation}

%%%%%%%% Modification paramétrage pour footnotes
% Le paquet semble déjà appelé en amont, curieux qu'il n'y ait pas eu de clash
% avec l'appel ci-dessous mais seulement quand j'ai essayé de l'inclure avec l'option lualatex

% \usepackage{hyperref}
% Parametrages
\hypersetup{
    colorlinks=true,% On active la couleur pour les liens. Couleur par défaut rouge
    linkcolor=black,% On définit la couleur pour les liens internes
    % filecolor=magenta,% On définit la couleur pour les liens vers les fichiers locaux      
    urlcolor=black,% On définit la couleur pour les liens vers des sites web
    % pdftitle={Puissance Quatre},% On définit un titre pour le document pdf
    % pdfpagemode=FullScreen,% On fixe l'affichage par défaut à plein écran
    }
% Pour avoir les numérotations arabic y compris dans les environnements tcolorbox

\usepackage{multicol} 
\usepackage{calc}
\usepackage{enumitem}
\usepackage{graphicx}
\usepackage{tabularx}
\usepackage{tablvar}

\setlength{\parindent}{0mm}
\renewcommand{\arraystretch}{1.5}
\renewcommand{\labelenumi}{\textbf{\theenumi{}.}}
\renewcommand{\labelenumii}{\textbf{\theenumii{}.}}
\setlength{\fboxsep}{3mm}

%\setlength{\headheight}{14.5pt}

\setlength{\premulticols}{0pt}

\newlength{\parindentMA}%
\setlength{\parindentMA}{\parindent}
\addtolength{\parindentMA}{30pt}

\usepackage{fancyhdr}
\fancyhead[C]{}
\fancyfoot{}

\newtheorem{exercice}{Exercice}

%%% MISE EN PAGE %%%
\usepackage{setspace}
\usepackage{pgf}%
\usepackage{pgfplots}
\pgfplotsset{compat=1.18}
\usetikzlibrary{babel,arrows, arrows.meta,calc,fit,patterns,plotmarks,shapes.geometric,shapes.misc,shapes.symbols,shapes.arrows,shapes.callouts, shapes.multipart, shapes.gates.logic.US,shapes.gates.logic.IEC, er, automata,backgrounds,chains,topaths,trees,petri,mindmap,matrix, calendar, folding,fadings,through,positioning,scopes,decorations.fractals,decorations.shapes,decorations.text,decorations.pathmorphing,decorations.pathreplacing,decorations.footprints,decorations.markings,shadows}

\usepackage[french]{babel}
\usepackage{icomma}
% loadPackagesFromContent

\newcounter{customfootnote}
\newcounter{tempreset}%
\newcounter{temp}

\newcommand{\e}{\mathrm{e}}
\usepackage{amsmath}
\usepackage{etoolbox}

%\usepackage{unicode-math}
%\usepackage{fontspec}
%\setmainfont{TeX Gyre Schola}
%\setsansfont{TeX Gyre Schola}
%\setmathfont{TeX Gyre Schola Math}
%\setmainfont{OpenDyslexic}[Scale=1.0]
%\setsansfont{OpenDyslexic}[Scale=1.0]
%\setmathfont{OpenDyslexic}[Scale=1.0,range=up/{Latin,latin,num}]

\begin{document}


{
  \normalsize
  \textbf{Seconde} \hfill \textbf{Fiche calcul 3} \hfill \textbf{Factorisation}
  \newline
}

\hrule

% @see : https://coopmaths.fr/alea?uuid=2e5df&id=2L11-5&n=3&d=10&s=4&alea=dOxp&cd=1&cols=2
\begin{exercice}
Factoriser les expressions suivantes.    
\begin{tasks}(3)
	\task  $A=-55n+88z$ 
	\task  $B=20n-35k$ 
	\task  $C=-8z^2+9z$ 
\end{tasks}
\tcbset{mathaleaexo/.style={breakable}}
% @see : https://coopmaths.fr/alea?uuid=0fe88&id=2L11-6&n=7&d=10&s=9&alea=g33Y&cd=1&cols=1
\end{exercice}

\begin{exercice}{}{}
Factoriser au maximum les expressions suivantes.


\begin{tasks}(2)
	\task $A=-2x^{2}-5x$
	\task $B=x(2x+5)-6x$
	\task $C=x^2-25$
	\task $D=(3x-5)(-2x+7)-(3x-5)(2x-1)$
	\task $E=(-x-9)(-5x-1)+(-x-9)(-4x-4)$
	\task $F=(-x-2)^2-(-x-2)(2x+7)$
\end{tasks}
\end{exercice}

\begin{exercice}
Utiliser $(a+b)^2 = a^2+2ab+b^2$ pour développer les expressions suivantes :
 \begin{tasks}(3)
    \task $(x+3)^2$
    \task $(4+a)^2$
    \task $(x+2y)^2$
    \task $(-2x+3y)^2$
    \task $\left(\dfrac{1}{3}+x\right)^2$
    \task $\left(\dfrac{x}{2}+\dfrac{y}{3}\right)^2$
 \end{tasks} 
\end{exercice}
 
\begin{exercice}
Utiliser $(a-b)^2 = a^2-2ab+b^2$ pour développer les expressions suivantes :
 
 \begin{tasks}(3)
    \task $(x-2)^2$
    \task $(1-a)^2$
    \task $(x-5y)^2$
    \task $(2x-3y)^2$
    \task $\left(\dfrac{1}{4}-x\right)^2$
    \task $\left(\dfrac{x}{2}-\dfrac{y}{3}\right)^2$
 \end{tasks} 
\end{exercice}
 
\begin{exercice}
Utiliser $(a+b)(a-b) = a^2-b^2$ pour développer les expressions suivantes :
 
 \begin{tasks}(3)
    \task $(x+7)(x-7)$
    \task $(x-4)(x+4)$
    \task $(x+2y)(x-2y)$
    \task $(2x+3y)(2x-3y)$
    \task $\left(\dfrac{3}{4}+x\right)\left(\dfrac{3}{4}-x\right)$
    \task $\left(\dfrac{3x}{2}+\dfrac{y}{5}\right)\left(\dfrac{3x}{2}-\dfrac{y}{5}\right)$
 \end{tasks} 
\end{exercice}
 
\begin{exercice}  
Factoriser les expressions suivantes.
\begin{tasks}(2)
    \task $A = 42z-49z^2$
    \task $B = 14c+63a$
\end{tasks}
\end{exercice}
 
\begin{exercice}  
Factoriser au maximum les expressions suivantes.
\begin{tasks}(2)
    \task $A = x(-6x-6)-7x$
    \task $B = (-x+6)(-4x-7)+(-x+6)(x+4)$
\end{tasks}
\end{exercice}
 
\begin{exercice}
Utiliser les identités remarquables pour factoriser les expressions suivantes :
 
 \begin{tasks}(3)
    \task $x^2+10x+25$
    \task $x^2+8x+16$
    \task $x^2-6x+9$
    \task $x^2-20x+100$
    \task $x^2-36$
    \task $x^2-\dfrac{1}{4}$
    \task $x^2-x+\dfrac{1}{4}$
    \task $9x^2+6x+1$
    \task $16x^2-\dfrac{4}{9}$
 \end{tasks} 
\end{exercice}

\begin{exercice}  
On considère l'expression littérale : $A(x) = x^2-6x-16$.
\begin{enumerate}
    \item Montrer que $A(x) = (x-3)^2-25$.
    \item Montrer que $A(x) = (x+2)(x-8)$.
    \item Utiliser l'expression de $A(x)$ la plus adaptée pour calculer $A(3)$, $A(8)$, $A(-2)$ et $A(0)$.
\end{enumerate}
\end{exercice}
 

\end{document}

% Local Variables:
% TeX-engine: luatex
% End: